\documentclass[11pt]{article}
\usepackage[letterpaper,margin=0.85in]{geometry}\usepackage[T1]{fontenc}\usepackage[utf8]{inputenc}\usepackage{enumitem,booktabs,array,xcolor,fancyhdr,tabularx}
\setlist{itemsep=2pt,topsep=3pt}\pagestyle{fancy}\fancyhf{}\renewcommand{\headrulewidth}{0.3pt}
\fancyhead[L]{\small Net Force and Acceleration (template)}\fancyhead[R]{\small Teacher key}\fancyfoot[C]{\small\thepage}
\definecolor{ans}{RGB}{28,125,74}\newcommand{\A}[1]{\textcolor{ans}{\textbf{#1}}}
\begin{document}

\section*{Net Force and Acceleration (template) --- teacher key}
\textbf{Question:} What decides how quickly something speeds up?\par
\subsection*{Outline and pacing (42 min)}\begin{tabular}{@{}llll@{}}\toprule Tab & Priority & Minutes & If behind\\\midrule
start & 1 & 2 & never trimmed\\
warmup & 1 & 4 & never trimmed\\
predict & 1 & 4 & never trimmed\\
explore & 3 & 8 & skipped first\\
build & 2 & 8 & shrinks to 2 of 4\\
practice & 1 & 8 & never trimmed\\
check & 1 & 4 & never trimmed\\
Huddle h1 & clock & minute 18 & pulled in by the clock\\
\bottomrule\end{tabular}
\subsection*{Targets}\begin{itemize}\item I can say what net force means.
\item I can predict how force and mass change acceleration.
\item I can calculate acceleration with a = F / m.\end{itemize}
\subsection*{Predictions}\begin{itemize}
\item A cart is pushed with a steady force. Now the force doubles. What happens to its acceleration? \A{It doubles.} a = F / m: double F, double a (mass is unchanged).
\item The same force pushes a cart. Now the cart carries twice the mass. What happens to its acceleration? \A{It halves.} a = F / m: double m, half a (force is unchanged).
\end{itemize}
\section*{Huddle script (answers in green)}
\subsection*{Definitions}\begin{itemize}
\item \textbf{Net force:} The \A{sum} of all the forces on an object.
\item \textbf{Acceleration:} How fast the \A{velocity} changes.
\end{itemize}
\subsection*{Worked example}
A 2 kg cart is pushed with a net force of 6 N. Find its acceleration.\\
\begin{itemize}
\item a = F / m
\item a = 6 / 2 = \A{3} m/s\textsuperscript{2}
\item In 1 s its speed rises by \A{3} m/s.
\end{itemize}
\subsection*{Label the cart}
\begin{enumerate}
\item The \A{applied force} pushes to the right.
\item \A{Friction} pushes to the left.
\end{enumerate}
\section*{Every level, with answers}
\subsection*{Find the acceleration}
\noindent\textbf{Level 1.} A 3 kg cart is pushed with a net force of 9 N. What is its acceleration?\\ \A{Answer: 3 m/s\textsuperscript{2}} \quad a = F / m = 9 / 3 = 3 m/s\textsuperscript{2}.\par\medskip
\noindent\textbf{Level 2.} A 0.5 kg cart feels a net force of 3 N. What is its acceleration?\\ \A{Answer: 6 m/s\textsuperscript{2}} \quad a = 3 / 0.5 = 6 m/s\textsuperscript{2}.\par\medskip
\noindent\textbf{Level 3.} A 5 kg cart is pushed right with 15 N while 5 N pushes left. What is its acceleration?\\ \A{Answer: 2 m/s\textsuperscript{2}} \quad Net force = 15 -- 5 = 10 N. a = 10 / 5 = 2 m/s\textsuperscript{2}.\par\medskip
\noindent\textbf{Level 4.} A rope pulls a 6 kg crate with 21 N. Friction is 3 N. What is the acceleration?\\ \A{Answer: 3 m/s\textsuperscript{2}} \quad Net = 21 -- 3 = 18 N. a = 18 / 6 = 3 m/s\textsuperscript{2}.\par\medskip
\subsection*{Find the force}
\noindent\textbf{Level 1.} A 3 kg cart accelerates at 3 m/s\textsuperscript{2}. What net force acts on it?\\ \A{Answer: 9 N} \quad F = m a = 3 $\times$ 3 = 9 N.\par\medskip
\noindent\textbf{Level 2.} A 0.5 kg cart accelerates at 4 m/s\textsuperscript{2}. What net force acts on it?\\ \A{Answer: 2 N} \quad F = 0.5 $\times$ 4 = 2 N.\par\medskip
\noindent\textbf{Level 3.} A 2500 g cart accelerates at 2 m/s\textsuperscript{2}. What net force acts on it, in newtons?\\ \A{Answer: 5 N} \quad 2500 g = 2.5 kg. F = 2.5 $\times$ 2 = 5 N.\par\medskip
\noindent\textbf{Level 4.} A 4 kg cart starts from rest and reaches 6 m/s in 2 s. What net force acted on it?\\ \A{Answer: 12 N} \quad a = 6 / 2 = 3 m/s\textsuperscript{2}. F = 4 $\times$ 3 = 12 N.\par\medskip
\section*{Common misconceptions}\begin{tabularx}{\textwidth}{@{}X X@{}}\toprule Students say & Say instead\\\midrule
A bigger force means a bigger speed. & Force sets how fast the speed changes, not the speed itself.\\
The heavier cart falls or moves faster. & More mass means less acceleration for the same force.\\
Subtract mass from force. & Divide. Check the units: N / kg = m/s\textsuperscript{2}.\\
\bottomrule\end{tabularx}
\end{document}